\Article{3}
Whenever the quantities \(v,v',v''\) etc. are completely independent of one another, the unknown quantity can be determined from them in only one way; consequently, in that case the uncertainty can neither be avoided nor diminished in any way, and nothing is left to arbitrary choice concerning the value of the unknown quantity to be deduced from the observations.

But the situation is very different whenever there is a mutual dependence among the quantities \(v,v',v''\) etc., which we shall assume to be expressed by \(\sigma\) conditional equations
\[
X=0,\quad Y=0,\quad Z=0\ \text{etc.}
\]
where \(X,Y,Z\) etc. denote given functions of the indeterminates \(v,v',v''\) etc. In this case our unknown quantity can be determined in infinitely many different ways by combinations of the quantities \(v,v',v''\) etc., since clearly, in place of the function \(u\), any other function \(U\) may be taken such that \(U-u\) vanishes identically when \(X=0,\ Y=0,\ Z=0\) etc. are set.

In applying this to a particular case, no difference in the value of the unknown quantity would result if the observations had absolute precision; but insofar as they remain subject to errors, a different combination will clearly in general give a different value of the unknown quantity. For example, in place of the error
\[
le+l'e'+l''e''+\text{etc.}
\]
which the function \(u\) had produced, we shall now have
\[
Le+L'e'+L''e''+\text{etc.}
\]
if we take the function \(U\) and denote the values of the derivatives \(\frac{\mathrm dU}{\mathrm dv},\frac{\mathrm dU}{\mathrm dv'},\frac{\mathrm dU}{\mathrm dv''}\) etc. respectively by \(L,L',L''\) etc. And although we cannot specify the errors themselves, it will nevertheless be possible to compare the mean errors to be fea%
\SourceLeaf{065}%
red in the different combinations of observations; the best combination will be that in which this mean error becomes as small as possible. Since it becomes
\[
=\BareRoot{(LLmm+L'L'm'm'+L''L''m''m''+\text{etc.})}
\]
the aim will be to make the sum \(LLmm+L'L'm'm'+L''L''m''m''+\text{etc.}\) attain its smallest value.

\Article{4}
Since the infinite variety of functions \(U\) which, under the condition stated in the preceding article, can take the place of \(u\) need be considered here only insofar as they give rise to different systems of values of the coefficients \(L,L',L''\) etc., it will first of all be necessary to investigate the relation that must hold among all admissible systems. Let us denote the fixed values of the partial derivatives
\begingroup\renewcommand{\arraystretch}{2}
\[
\begin{array}{ccc l}
\dfrac{\mathrm dX}{\mathrm dv},&\dfrac{\mathrm dX}{\mathrm dv'},&\dfrac{\mathrm dX}{\mathrm dv''}&\text{etc.}\\[5pt]
\dfrac{\mathrm dY}{\mathrm dv},&\dfrac{\mathrm dY}{\mathrm dv'},&\dfrac{\mathrm dY}{\mathrm dv''}&\text{etc.}\\[5pt]
\dfrac{\mathrm dZ}{\mathrm dv},&\dfrac{\mathrm dZ}{\mathrm dv'},&\dfrac{\mathrm dZ}{\mathrm dv''}&\text{etc. etc.}
\end{array}
\]
\endgroup
which they take when the true values are assigned to \(v,v',v''\) etc., respectively by
\[
\begin{array}{ccc l}
a,&a',&a''&\text{etc.}\\
b,&b',&b''&\text{etc.}\\
c,&c',&c''&\text{etc. etc.}
\end{array}
\]
It is clear that if the quantities \(v,v',v''\) etc. are conceived to receive increments \(\mathrm dv,\mathrm dv',\mathrm dv''\) etc. which leave \(X,Y,Z\) etc. unchanged, so that each remains \(=0\), i.e. increments satisfying the equations
\[
\begin{aligned}
0&=a\mathrm dv+a'\mathrm dv'+a''\mathrm dv''+\text{etc.}\\
0&=b\mathrm dv+b'\mathrm dv'+b''\mathrm dv''+\text{etc.}\\
0&=c\mathrm dv+c'\mathrm dv'+c''\mathrm dv''+\text{etc. etc.}
\end{aligned}
\]
then \(u-U\) must also remain unchanged, and therefore
\[
0=(l-L)\mathrm dv+(l'-L')\mathrm dv'+(l''-L'')\mathrm dv''+\text{etc.}
\]
\SourceLeaf{066}%
From this it is easily concluded that the coefficients \(L,L',L''\) etc. must be represented by formulas such as
\[
\begin{aligned}
L&=l+ax+by+cz+\text{etc.}\\
L'&=l'+a'x+b'y+c'z+\text{etc.}\\
L''&=l''+a''x+b''y+c''z+\text{etc.}
\end{aligned}
\]
etc., where \(x,y,z\) etc. denote fixed multipliers. Conversely, it is clear that if a system of fixed multipliers \(x,y,z\) etc. is taken arbitrarily, a function \(U\) can always be specified whose corresponding values of \(L,L',L''\) etc. conform to these equations, and which can take the place of \(u\) under the condition stated in the preceding article; indeed, this can be done in infinitely many different ways. The simplest way will be to set \(U=u+xX+yY+zZ+\text{etc.}\); more generally, one may set \(U=u+xX+yY+zZ+\text{etc.}+u'\), where \(u'\) denotes a function of the indeterminates \(v,v',v''\) etc. which always vanishes for \(X=0,\ Y=0,\ Z=0\) etc., and whose value in the particular case under consideration becomes a maximum or minimum. But for our purpose this makes no difference.

\clearpage
